% Preface for the First Workshop on Multilingual Report Generation via
% Retrieval Augmented Generation (RAG4Reports), co-located with ACL 2026.

Welcome to the First Workshop on Multilingual Report Generation via
Retrieval Augmented Generation (RAG4Reports), a half-day workshop
co-located with ACL 2026 in San Diego, CA, USA.

\begin{quote}
\textit{A workshop on multilingual long-form report generation focusing on
evaluating faithfulness and information coverage, tying together retrieval,
generation, and evaluation research.}
\end{quote}

% ---------------------------------------------------------------------------
% Motivation: why this workshop, why now
% ---------------------------------------------------------------------------
Incorporating external source information with parametric knowledge from large
language models to provide a comprehensive and grounded response to the users has
become a central component in modern AI applications.
\emph{Report generation} is a long-form
retrieval-augmented generation (RAG) task with strict attestation requirements
that makes it well-suited to explore questions of RAG evaluation and generation:
a long-form report summarizing the relevant information in a corpus is produced in
response to a report request, and the generated report should provide
proper attribution to source documents to establish trust.

RAG4Reports 2026 focuses on two urgent problems in report generation that
require a community effort to tackle: \emph{evaluation} and
\emph{multilinguality}. While numerous research papers on RAG have been
published in recent years, including proposed evaluation approaches, the
community has yet to reach a consensus on benchmarks and evaluation
measures for long-form, citation-grounded outputs. At the same time, as
generation models are increasingly multilingual and capable of incorporating
source information in different languages, report generation systems
should be fair to different languages and produce reports that the user can consume 
from relevant sources regardless of languages. The aim of this workshop is to bring
together researchers from information retrieval, natural language
processing, and applied domains to establish common ground for developing
and evaluating multilingual report generation systems.

% ---------------------------------------------------------------------------
% Submission statistics
% ---------------------------------------------------------------------------
The workshop solicited contributions through two tracks. The \emph{research
track} received 11 submissions through OpenReview, of which 9 were
accepted; of these, one was subsequently withdrawn and two were accepted as
non-archival submissions. The \emph{shared task track} attracted 21
registered teams, from which we received 8 system description papers. The
proceedings therefore contain 14 papers (6 research and 8 shared task), and
16 accepted papers are presented at the workshop as posters, with a subset
selected for oral presentation.

% ---------------------------------------------------------------------------
% Shared task overview
% ---------------------------------------------------------------------------
A central component of this workshop is the RAG4Reports shared task, which
builds on the 2024 NeuCLIR\footnote{\url{https://neuclir.github.io/}} and 
2025 RAGTIME\footnote{\url{https://trec-ragtime.github.io/}} shared tasks at the Text
Retrieval Conference (TREC) and extends them to the natural language
understanding community, with an explicit focus on non-English languages.
The shared task was hosted on TIRA\footnote{\url{https://www.tira.io/}},
which enables reproducible system evaluation through containerized
submissions.

The shared task consists of two sub-tasks. \emph{Task~A: Automatic Report
Evaluation} asks participants to rank system-generated reports from the
2025 TREC RAGTIME submissions against human annotations, accepting both
fully automatic and semi-automatic submissions (the latter using
organizer-provided essential facts). \emph{Task~B: Multilingual Report
Generation} asks participants to generate long-form, citation-grounded
reports from a corpus of four million English, Chinese, Russian, and Arabic
documents sampled from Common Crawl News (2021--2024), where each generated
sentence must be attributable to a cited source and the report must match
the language of the request. Submissions to Task~B were scored
automatically with the ARGUE framework and
its automatic implementation, Auto-ARGUE\footnote{\url{https://github.com/hltcoe/auto-argue}}, 
which combines nugget coverage
(whether information units expected of a good report are present) with
sentence support (whether each generated sentence is faithful to its cited
evidence).

% ---------------------------------------------------------------------------
% Keynote
% ---------------------------------------------------------------------------
We are delighted to host a keynote talk by Professor Chris Callison-Burch (University
of Pennsylvania) on the evaluation of long form RAG output, which speaks to the central 
theme of the RAG4Reports Workshop. Prof Callison-Burch has more than 200 
publications, which have been cited over 35,000 times. He is a Sloan Research Fellow, 
and he has received faculty research awards from Google, Microsoft, Amazon, Facebook, 
and Roblox, in addition to funding from DARPA, IARPA, and the NSF.

% ---------------------------------------------------------------------------
% Acknowledgments
% ---------------------------------------------------------------------------
This workshop would not have been possible without the contributions of many
people. We thank the program committee and external reviewers for their
careful and timely feedback; the shared task participants for engaging with
the benchmark and pushing its limits; and the TIRA team, particularly Maik Fröbe,
for hosting and supporting our evaluation infrastructure. 
We thank the ACL 2026 workshop chairs and
publication chairs for their guidance throughout the organization process. 
Finally, we thank all of the authors and shared task participants who submitted
their work. 

\bigskip

\noindent
We hope you enjoy the workshop and the discussions it sparks.

\bigskip

\noindent
\textit{The RAG4Reports 2026 Workshop Organizers}\\
Eugene Yang, Dawn Lawrie, Sean MacAvaney, James Mayfield,\\
Luca Soldaini, and Andrew Yates
